% !TEX program = xelatex
\documentclass[11pt,a4paper]{article}
\usepackage{algo-preamble}

\title{\textbf{L12 \textperiodcentered{} Άπληστοι Αλγόριθμοι II}\\[2pt]
\large Ελάχιστη Καθυστέρηση \& Τοπολογική Διάταξη}
\author{Algorithms Class Hub}
\date{\small Αλγόριθμοι και Πολυπλοκότητα (K17) \textperiodcentered{} ΕΚΠΑ}

\begin{document}
\maketitle

\begin{center}\itshape\small
Χρονοπρογραμματισμός με ελάχιστη μέγιστη καθυστέρηση (Earliest Deadline First), η
τεχνική του επιχειρήματος ανταλλαγής με αντιστροφές, οι τρεις τεχνικές απόδειξης
άπληστων αλγορίθμων, και η τοπολογική διάταξη σε DAG.
\end{center}

\section{Τι θα δούμε}

Δύο θέματα --- και μία τεχνική απόδειξης που θα ξαναδείς σε \textbf{κάθε}
εξεταστική:
\begin{enumerate}
  \item \textbf{Χρονοπρογραμματισμός με ελάχιστη καθυστέρηση} --- και το
  \textbf{επιχείρημα ανταλλαγής} (exchange argument), η δεύτερη μεγάλη τεχνική
  για άπληστους αλγορίθμους μετά το «greedy stays ahead» του προηγούμενου
  μαθήματος.
  \item \textbf{Τοπολογική διάταξη} --- πώς βάζεις σε σειρά εργασίες με
  προαπαιτούμενα, και γιατί αυτό είναι δυνατό \textbf{αν και μόνο αν} δεν
  υπάρχουν κύκλοι.
\end{enumerate}

\section{Χρονοπρογραμματισμός με Ελάχιστη Καθυστέρηση}

\subsection{Διατύπωση}

Μία μηχανή, που επεξεργάζεται \textbf{μία εργασία τη φορά}.
\begin{itemize}
  \item Η εργασία $j$ απαιτεί $t_j$ μονάδες χρόνου και έχει \textbf{προθεσμία}
  $d_j$.
  \item Αν ξεκινήσει τη στιγμή $s_j$, τελειώνει τη στιγμή $f_j = s_j + t_j$.
  \item \textbf{Καθυστέρηση} της $j$: $\ell_j = \max\{0,\ f_j - d_j\}$ --- πόσο
  \textbf{μετά} την προθεσμία της τελείωσε (ή $0$ αν πρόλαβε).
\end{itemize}
\textbf{Στόχος:} προγραμμάτισε όλες τις εργασίες ελαχιστοποιώντας τη
\textbf{μέγιστη} καθυστέρηση $L = \max_j \ell_j$.

\begin{intuition}
\textbf{Πρόσεξε τη διαφορά από το προηγούμενο μάθημα.} Εκεί διαλέγαμε ένα
υποσύνολο· εδώ προγραμματίζουμε \textbf{όλες} τις εργασίες --- απλώς αποφασίζουμε
τη \textbf{σειρά}. Και δεν μετράμε άθροισμα καθυστερήσεων αλλά τη
\textbf{χειρότερη} καθυστέρηση: μία πολύ αργοπορημένη εργασία χαλάει τη λύση,
ακόμη κι αν όλες οι άλλες είναι στην ώρα τους.
\end{intuition}

\subsection{Υποψήφια κριτήρια}

Πάλι το ίδιο άπληστο πρότυπο: βάλε τις εργασίες σε \textbf{κάποια σειρά}.

\begin{center}
\begin{tabular}{@{}l l l@{}}
\toprule
Κριτήριο & Σειρά & Αποτέλεσμα \\
\midrule
Μικρότερος χρόνος επεξεργασίας & αύξουσα $t_j$       & \textbf{αποτυγχάνει} \\
Μικρότερο χρονικό περιθώριο    & αύξουσα $d_j - t_j$ & \textbf{αποτυγχάνει} \\
Μικρότερη προθεσμία            & αύξουσα $d_j$       & \textbf{δουλεύει} $\checkmark$ \\
\bottomrule
\end{tabular}
\end{center}

\begin{warningbox}
\textbf{Γιατί αποτυγχάνει το «μικρότερος χρόνος επεξεργασίας»:} αγνοεί τελείως
την προθεσμία --- μια σύντομη εργασία με προθεσμία στο μέλλον προηγείται μιας
μακράς εργασίας που λήγει αμέσως. \textbf{Γιατί αποτυγχάνει το «μικρότερο
περιθώριο»:} μια εργασία με μικρό περιθώριο αλλά μακρινή προθεσμία μπορεί να
καθυστερήσει μία άλλη με κοντινή προθεσμία.
\end{warningbox}

\subsection{Ο αλγόριθμος: Earliest Deadline First}

\begin{keypoint}
\textbf{Άπληστος αλγόριθμος.} Ταξινόμησε τις εργασίες σε \textbf{αύξουσα σειρά
προθεσμίας} ($d_1 \le d_2 \le \dots \le d_n$) και προγραμμάτισέ τες \textbf{τη
μία μετά την άλλη}, χωρίς κενά.
\end{keypoint}

\begin{codebox}
\begin{verbatim}
ΕλάχιστηΚαθυστέρηση(t, d):
  Ταξινόμησε τις εργασίες ώστε d_1 <= d_2 <= ... <= d_n
  χρόνος <- 0
  Για j = 1 έως n:
    ανάθεσε στην εργασία j το διάστημα [χρόνος, χρόνος + t_j]
    s_j <- χρόνος ;  f_j <- χρόνος + t_j
    χρόνος <- χρόνος + t_j
  επίστρεψε τα διαστήματα [s_j, f_j]
\end{verbatim}
\end{codebox}

\begin{intuition}
\textbf{Διαισθητικά:} «κάνε πρώτα ό,τι λήγει πρώτο». Είναι ακριβώς αυτό που
κάνεις πριν την εξεταστική περίοδο --- το μάθημα με την κοντινότερη ημερομηνία
πάει πρώτο, ασχέτως πόσο διάβασμα θέλει.
\end{intuition}

\subsection{Δύο παρατηρήσεις-κλειδιά}

\begin{keypoint}
\textbf{Αδρανής χρόνος} (idle time): στιγμές που η μηχανή δεν δουλεύει.
\begin{itemize}
  \item Υπάρχει \textbf{πάντα} μια βέλτιστη λύση \textbf{χωρίς} αδρανή χρόνο (αν
  η μηχανή κάθεται, σπρώξε τα πάντα νωρίτερα --- καμία καθυστέρηση δεν
  μεγαλώνει).
  \item Ο \textbf{άπληστος} αλγόριθμος, εξ ορισμού, δεν έχει αδρανή χρόνο.
\end{itemize}
\end{keypoint}

\subsection{Αντιστροφές}

\begin{keypoint}
\textbf{Αντιστροφή} (inversion): ένα ζεύγος εργασιών $i, j$ με $d_i < d_j$ ---
δηλαδή η $i$ λήγει νωρίτερα --- αλλά η $j$ προγραμματίζεται \textbf{πριν} την $i$.
\end{keypoint}

Ο άπληστος αλγόριθμος, που ταξινομεί ακριβώς κατά προθεσμία, \textbf{δεν έχει
καμία αντιστροφή}. Αυτό είναι το «δαχτυλικό αποτύπωμα» της λύσης του.

\begin{intuition}
\textbf{Κρίσιμη παρατήρηση:} αν ένα χρονοδιάγραμμα χωρίς αδρανή χρόνο έχει έστω
\textbf{μία} αντιστροφή, τότε έχει και μία αντιστροφή με τις δύο εργασίες
\textbf{διαδοχικά} προγραμματισμένες. (Σκέψου το πιο «κοντινό» αντεστραμμένο
ζεύγος --- δεν μπορεί να έχει τίποτα ανάμεσά τους εκτός σειράς.)
\end{intuition}

\subsection{Το επιχείρημα ανταλλαγής}

\begin{keypoint}
\textbf{Ισχυρισμός.} Η αντιμετάθεση δύο \textbf{διαδοχικών, αντεστραμμένων}
εργασιών μειώνει τον αριθμό των αντιστροφών κατά $1$ και \textbf{δεν αυξάνει} τη
μέγιστη καθυστέρηση.
\end{keypoint}

\textbf{Απόδειξη.} Έστω $i, j$ διαδοχικές με αντιστροφή ($d_i < d_j$, αλλά η $j$
πριν την $i$). Αντιμεταθέτουμέ τες. Έστω $\ell_k$ οι παλιές και $\ell'_k$ οι
νέες καθυστερήσεις.
\begin{itemize}
  \item Για κάθε $k \ne i, j$: η εργασία $k$ δεν μετακινήθηκε καθόλου, άρα
  $\ell'_k = \ell_k$.
  \item Η $i$ μετακινήθηκε \textbf{νωρίτερα}, άρα τελειώνει νωρίτερα:
  $\ell'_i \le \ell_i$.
  \item Η $j$ μετακινήθηκε αργότερα. Τελειώνει τώρα εκεί όπου τελείωνε
  \textbf{πριν} η $i$ --- δηλαδή $f'_j = f_i$. Άρα:
  \[
  \ell'_j = f'_j - d_j = f_i - d_j \ \le\ f_i - d_i = \ell_i
  \]
  (η ανισότητα ισχύει αφού $d_j > d_i$).
\end{itemize}
Άρα και οι δύο νέες καθυστερήσεις, $\ell'_i$ και $\ell'_j$, είναι
$\le \ell_i \le$ της παλιάς μέγιστης. Η μέγιστη καθυστέρηση \textbf{δεν
αυξήθηκε}. $\blacksquare$

\begin{keypoint}
\textbf{Θεώρημα.} Η λύση $S$ του άπληστου αλγορίθμου είναι \textbf{βέλτιστη}.
\end{keypoint}

\textbf{Απόδειξη (εις άτοπον).} Έστω $S^*$ μια βέλτιστη λύση \textbf{χωρίς αδρανή
χρόνο} και με τον \textbf{ελάχιστο δυνατό αριθμό αντιστροφών}.
\begin{itemize}
  \item Αν το $S^*$ δεν έχει αντιστροφές: τότε είναι ταξινομημένο κατά
  προθεσμία, άρα $S = S^*$ --- ο άπληστος είναι βέλτιστος. $\checkmark$
  \item Αν το $S^*$ έχει τουλάχιστον μία αντιστροφή: τότε έχει και μία
  \textbf{διαδοχική} αντιστροφή. Αντιμεταθέτοντάς την, παίρνουμε λύση με
  \textbf{αυστηρά λιγότερες} αντιστροφές που \textbf{δεν} αυξάνει τη μέγιστη
  καθυστέρηση --- δηλαδή εξίσου βέλτιστη. Αυτό \textbf{αντιφάσκει} με την επιλογή
  του $S^*$ ως βέλτιστης με τις ελάχιστες αντιστροφές. $\blacksquare$
\end{itemize}

\subsection{Οι τρεις τεχνικές απόδειξης άπληστων αλγορίθμων}

\begin{keypoint}
\begin{enumerate}
  \item \textbf{Ο άπληστος προηγείται} (greedy stays ahead) --- δείχνεις ότι
  μετά από κάθε βήμα η μερική λύση του άπληστου είναι τουλάχιστον τόσο καλή όσο
  κάθε άλλης. \textit{(Προηγούμενο μάθημα, χρονοπρογραμματισμός διαστημάτων.)}
  \item \textbf{Επιχείρημα ανταλλαγής} (exchange argument) --- μετατρέπεις
  σταδιακά μια οποιαδήποτε βέλτιστη λύση σε αυτήν του άπληστου, χωρίς να χαλάει η
  βελτιστότητα. \textit{(Αυτή η διάλεξη.)}
  \item \textbf{Δομική εξήγηση} (structural bound) --- αποδεικνύεις ένα
  κάτω/άνω φράγμα μέσα στο πρόβλημα και δείχνεις ότι ο αλγόριθμος το πετυχαίνει.
  \textit{(Προηγούμενο μάθημα, διαμέριση διαστημάτων --- το βάθος.)}
\end{enumerate}
\end{keypoint}

\section{Τοπολογική Διάταξη}

\subsection{Το πρόβλημα}

Πολλές εργασίες με \textbf{προαπαιτούμενα}: η εργασία $u$ πρέπει να γίνει
\textbf{πριν} την $v$. Τα μοντελοποιούμε με ένα \textbf{κατευθυνόμενο γράφημα}
$G = (V, A)$: η ακμή $(u, v)$ σημαίνει «η $u$ προηγείται της $v$».

\begin{keypoint}
\textbf{Τοπολογική διάταξη} του $G$: μια διάταξη των κορυφών
$v_1, v_2, \dots, v_n$ τέτοια ώστε \textbf{κάθε} ακμή $(v_i, v_j)$ να πηγαίνει
«προς τα εμπρός» --- δηλαδή $i < j$.
\end{keypoint}

\begin{intuition}
\textbf{Διαισθητικά:} βάζεις όλες τις εργασίες σε μία γραμμή έτσι ώστε κάθε βέλος
να δείχνει \textbf{δεξιά}. Αν τα καταφέρεις, η σειρά αυτή είναι ένα έγκυρο
πρόγραμμα εκτέλεσης: κάθε εργασία γίνεται αφού έχουν τελειώσει όλα τα
προαπαιτούμενά της.

\textbf{Εφαρμογές:} γράφημα προαπαιτούμενων μαθημάτων, σειρά μεταγλώττισης
αρχείων κώδικα, διοχέτευση (pipeline) υπολογιστικών εργασιών, υπολογισμός
παράστασης όπως $((a+b)\cdot(c+d)+e)\cdot(a+b)$.
\end{intuition}

\subsection{Πότε υπάρχει;}

Δεν έχουν \textbf{όλα} τα κατευθυνόμενα γραφήματα τοπολογική διάταξη. Αν
$a \to b \to c \to a$, ποια μπαίνει πρώτη; Καμία --- υπάρχει \textbf{κύκλος}.

\begin{keypoint}
\textbf{Θεώρημα.} Ένα κατευθυνόμενο γράφημα έχει τοπολογική διάταξη \textbf{αν
και μόνο αν} είναι \textbf{άκυκλο} (Directed Acyclic Graph --- DAG).
\end{keypoint}

\textbf{Μέρος 1 --- αν έχει τοπολογική διάταξη, τότε είναι DAG.} Έστω ότι το $G$
έχει τοπολογική διάταξη $v_1, \dots, v_n$ \textbf{και} έναν κύκλο $C$. Έστω
$v_i$ η κορυφή του $C$ με τον \textbf{μικρότερο} δείκτη, και $v_j$ η κορυφή
ακριβώς \textbf{πριν} την $v_i$ μέσα στον $C$. Τότε η $(v_j, v_i)$ είναι ακμή.
Από την επιλογή της $v_i$ ως μικρότερης, $i < j$. Αλλά αφού $(v_j, v_i)$ είναι
ακμή σε τοπολογική διάταξη, $j < i$. \textbf{Άτοπο.} $\blacksquare$

\textbf{Μέρος 2 --- κάθε DAG έχει κορυφή χωρίς εισερχόμενες ακμές.} Έστω ότι
\textbf{κάθε} κορυφή έχει εισερχόμενη ακμή. Ξεκίνα από μια κορυφή $v$ και
ακολούθησε ακμές \textbf{προς τα πίσω}: αφού η $v$ έχει εισερχόμενη $(u, v)$,
πήγαινε στην $u$· αφού η $u$ έχει εισερχόμενη, συνέχισε. Με $n$ κορυφές, μετά από
$n$ βήματα θα συναντήσεις κάποια κορυφή \textbf{δύο φορές} --- και το κομμάτι
ανάμεσα στις δύο επισκέψεις είναι \textbf{κύκλος}. Άτοπο, αφού το $G$ είναι DAG.
$\blacksquare$

\textbf{Μέρος 3 --- κάθε DAG έχει τοπολογική διάταξη.} Επαγωγή στο πλήθος
κορυφών $n$.
\begin{itemize}
  \item \textit{Βάση:} $n = 1$ --- προφανές.
  \item \textit{Επαγωγικό βήμα:} έστω DAG με $k+1$ κορυφές. Από το Μέρος 2,
  υπάρχει κορυφή $v$ \textbf{χωρίς εισερχόμενες} ακμές. Το $G \setminus \{v\}$
  είναι DAG με $k$ κορυφές, άρα (επαγωγική υπόθεση) έχει τοπολογική διάταξη.
  Τοποθετούμε την $v$ \textbf{πρώτη}: αφού δεν έχει εισερχόμενες ακμές, κάθε ακμή
  της δείχνει προς τα εμπρός. Η διάταξη είναι έγκυρη. $\blacksquare$
\end{itemize}

\subsection{Ο αλγόριθμος}

Η απόδειξη του Μέρους 3 \textbf{είναι} ο αλγόριθμος: βρίσκεις επανειλημμένα μια
κορυφή χωρίς εισερχόμενες ακμές και την «βγάζεις».

\begin{codebox}
\begin{verbatim}
ΤοπολογικήΔιάταξη(G):
  βρες κορυφή v του G χωρίς εισερχόμενες ακμές
  L <- ΤοπολογικήΔιάταξη(G \ {v})
  επίστρεψε v, L          // η v μπροστά, μετά η υπόλοιπη λίστα
\end{verbatim}
\end{codebox}

\begin{keypoint}
\textbf{Πολυπλοκότητα: $O(m + n)$.} Με αφελή υλοποίηση θα ήταν $O(n^2)$ (ψάχνεις
κάθε φορά από την αρχή). Το κόλπο: διατήρησε
\begin{itemize}
  \item $\text{count}[w] = $ πλήθος εναπομεινάντων \textbf{εισερχόμενων} ακμών
  της $w$·
  \item $S = $ σύνολο (λίστα) κορυφών με $\text{count} = 0$.
\end{itemize}
\textbf{Αρχικοποίηση:} μία σάρωση του γραφήματος, $O(m + n)$. \textbf{Διαγραφή
της $v$:} βγάλε την από το $S$ ($O(1)$)· για κάθε ακμή $(v, w)$ μείωσε το
$\text{count}[w]$, και αν φτάσει $0$ βάλε την $w$ στο $S$ --- $O(1)$ ανά ακμή.
Συνολικά $O(m + n)$.
\end{keypoint}

\begin{recapbox}
\textbf{Τι κρατάμε από αυτή τη διάλεξη}
\begin{enumerate}
  \item \textbf{Ελάχιστη μέγιστη καθυστέρηση} --- μία μηχανή, εργασίες με
  $(t_j, d_j)$, καθυστέρηση $\ell_j = \max\{0, f_j - d_j\}$. Σωστό κριτήριο:
  \textbf{Earliest Deadline First} (αύξουσα προθεσμία).
  \item \textbf{Αντιστροφή} $=$ ζεύγος εργασιών εκτός σειράς προθεσμίας. Ο
  άπληστος δεν έχει καμία.
  \item \textbf{Επιχείρημα ανταλλαγής} --- αντιμετάθεση μιας διαδοχικής
  αντιστροφής μειώνει τις αντιστροφές κατά 1 χωρίς να μεγαλώνει το $L$. Άρα κάθε
  βέλτιστη λύση μπορεί να γίνει η λύση του άπληστου.
  \item \textbf{Τρεις τεχνικές απόδειξης:} greedy stays ahead, exchange
  argument, structural bound.
  \item \textbf{Τοπολογική διάταξη} --- διάταξη κορυφών ώστε κάθε ακμή να δείχνει
  εμπρός. Υπάρχει \textbf{αν και μόνο αν} το γράφημα είναι \textbf{DAG} (άκυκλο).
  \item Αλγόριθμος: βγάζε επανειλημμένα κορυφή χωρίς εισερχόμενες ακμές ---
  $O(m + n)$ με \texttt{count[]} και σύνολο πηγών.
\end{enumerate}
\end{recapbox}

\lecturefooter
\end{document}
